\documentclass[11pt]{article}
\usepackage{amsmath,amssymb,amsthm,hyperref}
\newtheorem{lemma}{Lemma}
\newtheorem{proposition}{Proposition}
\begin{document}
\title{Erd\H{o}s Problem 869: a checked attempt}
\author{GPT\_6\_Astra\_Ultra}
\date{24 September 2026}
\maketitle

\section*{Statement and approach}
Must a union of two disjoint asymptotic bases of order two contain a minimal subbasis? No. We reconstruct the relevant case of Daniel Larsen's \emph{Three Questions of Erd\H{o}s--Nathanson on Asymptotic Bases of Order 2}, \url{https://raw.githubusercontent.com/Larsen-Daniel/Erdos-869/main/869.pdf}. The block construction and its probabilistic verification are included. Representations are unordered; a balanced pair has ratio at most $100$.

\section*{A flexible block construction}
Put $N_j=4^{j+1}$. Partition the positive integers independently and uniformly into $X_1,X_2,X_3$. At stage $j$ let $N=N_j$ and let $B_{<j},C_{<j}$ be the already constructed disjoint sets below $N$; write $A_{<j}=B_{<j}\cup C_{<j}$. Choose disjoint sets $G_j\subseteq B_{<j}$, $H_j\subseteq C_{<j}$ by any rule depending on the past. Set
\[
 V_j=(4N/3,2N)\cup(8N/3,3N)\cup
       \bigl(4N-((0,N)\setminus A_{<j})\bigr),
\]
and add $(V_j\cap X_1)\cup(4N-G_j)$ to $B$ and $(V_j\cap X_2)\cup(4N-H_j)$ to $C$. Intervals denote their integer points. The new additions lie strictly between $N$ and $4N$. They are disjoint: the planted complements have partners in $A_{<j}$, whereas the random upper additions have partners outside it. In particular $B\cap C=\varnothing$ throughout.

The complete representations of $4N$ in $B\cup C$ are exactly
\[
 \{f,4N-f\}\quad(f\in G_j\cup H_j).\tag{1}
\]
Indeed two old elements sum below $2N$. Among new lower elements, the first interval added to itself sums below $4N$, the first plus the second sums above $4N$, and two from the second are still larger. An old element and either lower interval sum below $4N$. Finally the upper interval's partner lies below $N$, and its membership rule leaves exactly the planted pairs. Future stages cannot affect this list.

We now prove that, for a suitable realization, both colors have at least $c n$ balanced representations of every sufficiently large $n\notin\{N_j\}$, with an absolute $c>0$, irrespective of the selection rule. At stage $j$ the set $B$ contains $X_1\cap I$ for
\[
 I_0=(N/24,3N/64),\ I_1=(2N/3,3N/4),\
 I_2=(4N/3,2N),\ I_3=(8N/3,3N).
\]
The first two come from stages $j-3,j-1$. It also contains
\[
 X_1\cap(4N-(X_3\cap J)),\qquad J=(N/4,3N/4).
\]
To see the latter, no earlier planted complement lies in $J$: the immediately previous ones lie in $(3N/4,N)$ and all older ones below $N/4$. Every random addition in $J$ lies in $X_1\cup X_2$, so $X_3\cap J$ is absent from $A_{<j}$.

For two of these ordinary intervals, each integer in the sum interval with one percent of its length trimmed from each end has $\gg N$ candidate unordered pairs, with an absolute constant (all interval lengths are at least $N/192$). Distinct candidate pairs use disjoint coordinates, apart from at most one diagonal candidate, which we discard. Their probability of being in $X_1$ is $1/9$, independently. For $I+(4N-J)$ and a sum $4N+\Delta\ne4N$, there are $\gg N$ candidates $y\in J$, $y+\Delta\in I$. The required colors are $y\in X_3$, $y+\Delta,4N-y\in X_1$. After discarding at most one candidate with coincident coordinates, a greedy selection leaves $\gg N$ triples whose coordinate sets are disjoint: each triple conflicts with at most eight other triples, because its three affine coordinates are injective functions of $y$. Their success probability is $1/27$ independently. Chernoff's inequality gives failure probability $\exp(-c_0N)$ for a fixed target in either case, and the same argument applies to $C$.

For completeness, the following open intervals, written in units of $N$, lie in the trimmed sum intervals just described:
\[
\begin{array}{c|c}
I_0+I_2&(1.39,2.04)\\
I_1+I_2&(2.01,2.74)\\
I_2+I_2&(2.69,3.98)\\
I_1+(4N-J)&(3.93,4.49)\setminus\{4\}\\
I_3+I_2&(4.02,4.98)\\
I_2+(4N-J)&(4.60,5.48)\\
I_3+I_3&(5.35,5.99)\\
I_3+(4N-J)&(5.93,6.70)\setminus\{4\}.
\end{array}
\]
They cover $[1.5,6]\setminus\{4\}$ with overlap. All summands are between $N/24$ and $15N/4$, so their ratio is at most $90<100$. The union bound over $O(N)$ integer targets and the summability of $N_j\exp(-c_0N_j)$ prove the assertion for all sufficiently large stages. The ranges $[3N_j/2,6N_j]$ meet consecutively and cover a tail. This establishes the claimed uniform density of balanced representations without assuming independence of the adaptively selected sets.

\section*{Selecting pairs to prevent minimality}
At stage $j$, call a pair $\{b,c\}$ eligible if $b\in B_{<j}$, $c\in C_{<j}$, and $\min(b,c)\ge j$. Retain every pair once it becomes eligible. Select a previously unselected eligible pair with smallest maximum (break ties deterministically), and take $G_j=\{b\}$, $H_j=\{c\}$. If none exists use empty sets. Eventually a pair always exists: the representation bound at $N_j-1$ implies that each old color has $\gg N_j$ elements, so even after deleting elements below $j$ it supplies far more than $j$ eligible pairs, while only $j-1$ pairs have been used. Every eligible pair is eventually selected, since only finitely many pairs have bounded maximum.

Consequently each color represents all large nonspecial integers by the block estimate, and each large $N_{j+1}$ by its planted pair. Thus $B,C$ are disjoint bases. Moreover the least element of the selected pair tends to infinity. For any fixed $T$, a pair containing an element at most $T$ must have become eligible by stage $T$, hence has maximum at most $N_T$; only finitely many such pairs exist.

Let $D\subseteq B\cup C$ be any subbasis. The set $D$ must meet every eligible pair whose maximum is sufficiently large, regardless of when it first became eligible. Indeed, when a pair $\{b,c\}$ is selected at stage $i$, its witness satisfies $N_{i+1}>4\max(b,c)$; if $D$ omitted both endpoints, this sufficiently large witness would have no representation, by (1). For a large integer $M$, take the least $j$ with $N_j>100M$. Since $j=O(\log M)<M$, every pair consisting of one element of each color in $[M,100M]$ is eligible at stage $j$. It follows that at least one of the two colors is retained in its entirety by $D$ on this interval.

For any large nonspecial $u$, put $M=\lceil u/101\rceil$. Every balanced representation of $u$ has both summands in $[M,100M]$: the smaller is at least $u/101$ and the larger at most $100u/101\le100M$. Therefore $D$ has $\gg u$ representations of $u$. Removing any fixed $d\in D$ loses at most one of them. For special integers $N_{j+1}$, all summands in (1) tend to infinity (the planted large summands exceed $3N_j$), so removing $d$ eventually loses none. Thus $D\setminus\{d\}$ remains a basis. No subbasis $D$ is minimal.

\section*{Assessment}
The counterexample is completely reconstructed from the attributed block strategy. The probability estimates use only independent finite-coordinate trials, Chernoff's inequality and a summable union bound. No novel-resolution claim is made.

\textbf{Completion rate estimate: 100\%.}

\end{document}
